\documentclass[11pt]{article}
\usepackage[margin=25mm]{geometry}
\usepackage{amsmath,amssymb,amsthm,mathtools,hyperref}
\newcommand{\C}{\mathbb C}
\newcommand{\Tr}{\operatorname{Tr}}
\newcommand{\im}{\operatorname{im}}
\newcommand{\cond}{\operatorname{cond}}
\newcommand{\lc}{\operatorname{lc}}
\newcommand{\norm}[1]{\left\lVert#1\right\rVert}
\title{Metric cost and the original marked observation}
\author{Split-Zero research programme}\date{19 September 2026}
\begin{document}\maketitle

This is the complete receiving proof for Sections 4--5 of the
supplied low-cutoff continuation \cite{source}. LCS1--40, ORE1--23 and the
independent conductor calculation supply its original source
objects and estimates. PT1--41 \cite{PT} supplies the full-primary
positive-trace minimum, including all nilpotent blocks.
The final selected-minor conclusion is corrected through the
exact COI11 map \cite{COI} instead of silently identifying two coefficient
sections. All observations below retain their original period.

\section{The cost of a positive-trace metric}
Use the full original quartet polynomial with fixed multiplicity,
$e=1+k(m_0-1)$, $q=e(k+1)^2$, $c=k/2$, and source cutoff $q-1$.
Let $G_0=G_{q-1}$ denote its original arithmetic minimum metric.
In its orthonormal coordinates the full multiplication operator
satisfies the proved original boundary identity NCE10:
\begin{equation}\tag{CMO1}
 \begin{gathered}
 B_0=G_0^{1/2}(A-cI)G_0^{-1/2}=iJ+\epsilon_0fe^*,\\
 J=J^*,\quad \norm e=\norm f=1,\quad e^*f=0,\quad
 \epsilon_0^2=\frac{\nu_0-\omega_q}{\omega_{q-1}} .
 \end{gathered}
\end{equation}
In this section the two occurrences of $e$ have different types:
the integer multiplicity in the first paragraph and the unit
boundary vector in CMO1. The integer is never used as a vector.
The exact source Gram and LCS9 give
\begin{equation}\tag{CMO2}
 \epsilon_0^2=\frac{1-d_0}{m_zp_{1,q-1}},\qquad
 \epsilon_0=\frac{4q}{\pi\sqrt{p_{1,q-1}}}(1+o_h(1)).
\end{equation}
All norms retain their source mass. The asymptotic follows from
$Y^2m_z\to1$, $Y=4q/\pi$, and $d_0\to0$, proved in LCS19,30.
The original multiplication bound gives
\[
 \norm J\le J_*=
 C_{\rm mult}(h)(2q-2+\lfloor k/3\rfloor).
\]

For any positive Hermitian metric $\widehat G$ on the same
original coefficient space set
\[
 M=G_0^{-1/2}\widehat G G_0^{-1/2},\qquad
 S=M^{1/2},\qquad \kappa=\cond M,\qquad
 B_S=SB_0S^{-1},\qquad H_S=B_S+B_S^* .
\]
Thus $S$ is the exact isometry from the new metric in the old
orthonormal coordinates to a Euclidean coordinate space.
The weight has no extra factor $1/2$.
Let $\mathcal P(\widehat G)=\Tr(H_S)_+$.
The spectral trace is invariant under this coordinate
identification, so it equals the positive trace for $\widehat G$.
The original symmetric root multiset has $\Re\Tr B_0=0$,
including every repeated primary; hence $\Tr H_S=0$.
If its real eigenvalues are $h_j$, this says
$\sum_{h_j>0}h_j=-\sum_{h_j<0}h_j$.
Every $|h_j|$ is bounded by that common sum, proving
$\norm{H_S}\le\mathcal P(\widehat G)$.

Put $u=Sf$, $w=S^{-*}e$. Then
$w^*u=e^*f=0$ and
\[
 SB_0S^{-1}=iSJS^{-1}+\epsilon_0uw^*.
\]
On the orthonormal pair $u/\norm u,w/\norm w$ the Hermitian
rank-two matrix $uw^*+wu^*$ has eigenvalues
$\pm\norm u\norm w$; it is zero on their orthogonal complement.
Moreover
$\norm u\norm w\ge\sigma_{\min}(S)/\sigma_{\max}(S)
=\kappa^{-1/2}$, while
$\norm{iSJS^{-1}+(iSJS^{-1})^*}\le2J_*\sqrt\kappa$.
The triangle inequality applied to this exact decomposition proves
\begin{equation}\tag{CMO3}
 \boxed{\epsilon_0\le
       \mathcal P(\widehat G)\sqrt\kappa+2J_*\kappa.}
\end{equation}
The classical Hermitian norm and quadratic-coordinate framework is
credited to Boyd and Vandenberghe \cite[Appendix A.5]{Boyd};
the entire rank-two comparison just used is explicitly proved here.

PT1--41 evaluates the full original minimum as
$L_{h,k}=\delta q(k+1)/2$. For every attaining metric, put
$x=\sqrt\kappa$. CMO3 gives $2J_*x^2+L_{h,k}x-\epsilon_0\ge0$.
Since $x>0$ and $J_*>0$, solving the quadratic with its positive
root and rationalizing proves the finite bound
\begin{equation}\tag{CMO4}
 \boxed{\kappa\ge
 \left(\frac{2\epsilon_0}
 {L_{h,k}+\sqrt{L_{h,k}^2+8J_*\epsilon_0}}\right)^2.}
\end{equation}
LCS19 gives
$\log\epsilon_0=q\log(4/\pi)+O_h(k+\log q)$.
Thus $L_{h,k}^2/(J_*\epsilon_0)\to0$ at every fixed
multiplicity, since $L_{h,k}$ and $J_*$ are polynomial in $k$.
CMO2 and $J_*/q\to2C_{\rm mult}(h)$ now give
\begin{equation}\tag{CMO5}
 \kappa\ge
 \frac{1-o_h(1)}{\pi C_{\rm mult}(h)\sqrt{p_{1,q-1}}},
 \qquad
 \boxed{\liminf_{k\to\infty}\frac{\log\kappa}{q}
                         \ge\log(4/\pi)>0.}
\end{equation}
The same conclusion follows for any family with
$\mathcal P(\widehat G)$ bounded by a polynomial in $q$:
replace $L_{h,k}$ by that actual value in the positive root.
This does not change the full multiplication operator or remove
nilpotents. It quantifies the isometry's distortion in its original
metric, while preserving the valid positive-trace minimum.

\section{Conductor detection on original low degrees}
Now retain the simple-packet conductor domain of FEX1--4:
$q=(k+1)^2$, $q'=(k-7)^2$, original $\Lambda_k:E\to B=\im\Lambda_k$,
$K=\ker\Lambda_k$, conductor
\[
 \mathcal T_AP(S')=\sum_{\nu}a_\nu P(S'+b_\nu),\qquad
 \mu_j=\sum_\nu a_\nu b_\nu^j,\qquad
 v=\min\{j:\mu_j\ne0\}\le80 .
\]
Here $v$ is the conductor order. Every representative of a class
in $K$ satisfies $\chi_{k-8}\mid\mathcal T_AP$.
This implication descends to $E$: adding $\chi_kH$ vanishes
at each shifted lower root, exactly as proved in the preceding
conductor source.
Define the literal conductor map
$C:E\to\C[S']/(\chi_{k-8})$ by $C[P]=[\mathcal T_AP]$.
Then $K\subset\ker C$, so there is a unique linear map
\begin{equation}\tag{CMO6}
 \bar C:B\longrightarrow\C[S']/(\chi_{k-8}),\qquad
 \bar C(\Lambda_kx)=Cx,\qquad C=\bar C\Lambda_k .
\end{equation}
Well-definedness follows because any two lifts differ by $K$;
surjectivity of $\Lambda_k$ onto $B$ proves uniqueness.
This is an exact typed factorization of the original maps.

If $P$ has degree $D\ge v$ and nonzero leading coefficient,
Taylor expansion gives
\[
 \deg\mathcal T_AP=D-v,\qquad
 \lc(\mathcal T_AP)=\mu_v\binom Dv\lc(P)\ne0 .
\]
For $v\le D<q'+v$ this degree is below $q'$, so
divisibility by $\chi_{k-8}$ is impossible.
If $D<v$, every derivative term carrying a nonzero moment
vanishes. Consequently, with $\mathcal P_{-1}=0$,
\begin{equation}\tag{CMO7}
 K\cap\mathcal P_D=K\cap\mathcal P_{v-1}
 \quad(v-1\le D<q'+v),\qquad
 CS^v=\mu_v[1]\ne0 .
\end{equation}
In particular
\begin{equation}\tag{CMO8}
 j_*=\min\{j:\Lambda_k S^j\ne0\}\le v\le80 .
\end{equation}
All polynomial spaces here are included by their actual classes
in $E$, so there is no quotient-metric identification in CMO7.

\section{The marked derivative and its exact observed coefficient}
Use the marked family in the original increasing-power frame:
\[
 A(t)=A+t e_0e_{q-1}^{\mathsf T},\qquad
 U'(t)=A(t)U(t)/u,\quad U(0)=I,\quad u\ne0,
 \qquad \Lambda_h(t)=\Lambda_kU(t)^{-1}.
\]
Writing $E=e_0e_{q-1}^{\mathsf T}$, the coefficients of
$U(t)=\sum_{n\ge0}U_nt^n$ satisfy $U_0=I$, $U_{-1}=0$ and
$(n+1)U_{n+1}=(AU_n+EU_{n-1})/u$.
Their norms are termwise bounded by the nonnegative coefficients
of $\exp((\norm A)t/|u|+(\norm E)t^2/(2|u|))$, since those
coefficients satisfy the corresponding scalar recurrence.
Thus the series converges for every $t$, solves the equation
term by term, and is unique by the same recurrence.
It has an analytic inverse near $0$, since its determinant
starts at one.
Let $v_\lambda=[b_\lambda\chi_k/(S-\lambda)]$ be the full
terminal eigenclass and put
$Y(t)=e^{\lambda t/u}\Lambda_h(t)v_\lambda-\Lambda_kv_\lambda$.
Differentiating the inverse on its actual domain gives
$\Lambda_h'=-\Lambda_h A(t)/u$.
Since $Av_\lambda=\lambda v_\lambda$ and
$e_{q-1}^{\mathsf T}v_\lambda=b_\lambda$, we obtain
\begin{equation}\tag{CMO9}
 Y'(t)=-\frac{b_\lambda}{u}\,t\,e^{\lambda t/u}
                               \Lambda_h(t)1 .
\end{equation}
On the admitted conductor domain $k\ge9$, we have
$q\ge100>v+1$ since $v\le80$.
For $j\le v$, the Taylor recurrence from
$\Lambda_h'=-\Lambda_h(A+t e_0e_{q-1}^{\mathsf T})/u$ yields
\begin{equation}\tag{CMO10}
 \Lambda_h^{(j)}(0)1=(-u)^{-j}\Lambda_kS^j.
\end{equation}
To justify the omitted marked terms explicitly, each term
containing $E=e_0e_{q-1}^{\mathsf T}$ and acting on $1=e_0$
has a rightmost $E$ preceded on its right by fewer than $q-1$
copies of $A$. But $A^r1=S^r$ for $r<q$ and
$e_{q-1}^{\mathsf T}S^r=0$ for $r<q-1$.
Every such term vanishes. The remaining term is
$(-u)^{-j}\Lambda_kA^j1$, proving CMO10 with its sign.

For $j<j_*$ all observed Taylor coefficients vanish.
Leibniz's rule at $j_*$ in CMO9 therefore gives the first
nonzero derivative
\begin{equation}\tag{CMO11}
 \boxed{Y^{(j_*+2)}(0)
 =b_\lambda(j_*+1)(-u)^{-(j_*+1)}\Lambda_kS^{j_*}\ne0,
 \qquad j_*+2\le v+2\le82.}
\end{equation}
Applying $\bar C$, the same reasoning uses $CS^j=0$ for
$j<v$ and $CS^v=\mu_v[1]$. It proves the explicit coefficient
\begin{equation}\tag{CMO12}
 \bar C Y^{(r+2)}(0)=0\quad(r<v),\qquad
 \boxed{\bar C Y^{(v+2)}(0)
 =b_\lambda(v+1)(-u)^{-(v+1)}\mu_v[1]\ne0.}
\end{equation}
For $v=0$ this also proves $\Lambda_k1\ne0$ on every
admitted conductor degree. No arithmetic projected sign is
inferred from this marked analytic Taylor coefficient.

\section{The precise return to the original coefficient image}
COI proves $\im L_0=\mathcal P_{m-1}$ for its consecutive
comparison section, $m=16k-48-v$. At $k\ge29$ and $0\le v\le80$,
$v-1\le m-1<q'+v$: the upper difference is
$q'+v-(m-1)=k^2-30k+98+2v$, positive at $29$
and increasing thereafter. CMO7 therefore proves
\begin{equation}\tag{CMO13}
 K\cap\im L_0=K\cap\mathcal P_{v-1},\qquad
 \operatorname{rank}(\Lambda_k L_0)=m-t_0,\qquad
 t_0=\dim(K\cap\mathcal P_{v-1}).
\end{equation}
For the original selected minor $I$, use the proved map
and its observed image, without dropping the conductor term:
\begin{equation}\tag{CMO14}
 \boxed{L_I=F_CD_{0I}+L_0M_{0I},\qquad
 \Lambda_kL_I=(\Lambda_kF_C)D_{0I}
                          +(\Lambda_kL_0)M_{0I}.}
\end{equation}
Here $M_{0I}$ is invertible by the full COI9--11 transition.
Thus the supplied note's unqualified consecutive-image statement
is replaced by its exact original-minor receiver.

Finally RBV's explicit polynomial \cite{RBV} is
\[
 t_0'=\frac{\pi}{2\gamma},\qquad
 P_k(S)=\sum_{j=0}^{4k}\frac{(t_0')^{2j}(S-c)^{2j}}{(2j)!}.
\]
It has degree $8k$, nonzero leading coefficient
$(t_0')^{8k}/(8k)!$, and the negative reflected trace proved
in the cited complete source. For $k\ge29$,
$v\le8k<q'+v$ and $8k\le m-1$, directly from the above
degree formulas. Its conductor has nonzero leading coefficient
\begin{equation}\tag{CMO15}
 \mu_v\binom{8k}{v}\frac{t_0'^{\,8k}}{(8k)!}\ne0 .
\end{equation}
Its degree is $8k-v<q'$, so $CP_k\ne0$ and hence
$\Lambda_kP_k\ne0$ by CMO6. The negative-trace test survives
the actual observation on this domain. Its signed residue
pairing and the arithmetic projected current retain their
respective exact definitions; CMO6 and CMO14 are the proved
observation maps, and do not assert equality of those pairings.

\begin{thebibliography}{9}
\bibitem{Boyd} Stephen Boyd and Lieven Vandenberghe,
\emph{Convex Optimization}, Cambridge University Press, 2004,
Appendix A.5,
\url{https://web.stanford.edu/~boyd/cvxbook/bv_cvxbook.pdf}.
\bibitem{source} Supplied programme continuation,
\emph{Low-cutoff energies, conductor and marked defect},
19 September 2026, equations (35)--(47).
The complete receiving proof and selected-minor correction
are given above.
\bibitem{PT} Split-Zero programme,
\emph{The full-primary positive-trace minimum and its original-metric
receivers}, PT1--41, 19 September 2026.
\bibitem{COI} Split-Zero programme,
\emph{Every original coefficient section}, COI1--17,
19 September 2026.
\bibitem{RBV} Split-Zero programme,
\emph{The original residue boundary value}, RBV1--36,
19 September 2026.
\end{thebibliography}
\end{document}
